\chapter{АРХИТЕКТУРА И РЕАЛИЗАЦИЯ ПЛАТФОРМЫ}

\section{Декомпозиция на модули}

Платформа разделена на три самостоятельных модуля, связанных явными контрактами:

\begin{itemize}
    \item \textbf{generation\_core} --- каноническое представление, LLM-провайдеры, построители, проверка, симуляция, импорт/экспорт, CLI и внутренний REST API;
    \item \textbf{backend} --- email/OTP-аутентификация, пользователи, беседы, сообщения, версии моделей, симуляции и вызовы generation\_core на FastAPI \cite{fastapi2024};
    \item \textbf{frontend} --- интерфейс бесед, визуализация CLD/SFD, форматирование уравнений, запуск симуляций и экспорт результатов.
\end{itemize}

Generation core не зависит от пользователей, JWT и продуктовых сущностей backend. Это позволяет запускать его отдельно в CLI, тестировать как вычислительный модуль и заменять LLM-провайдера без изменения интерфейса пользователя. Backend вызывает его как внутренний сервис, а не как источник кода, загружаемого пользователем.

PostgreSQL \cite{postgres2024} хранит пользователей и артефакты, Redis используется для очередей и состояния задач, а Taskiq \cite{taskiq2024} выполняет длительные операции асинхронно. Разделение вычислительного контура и контура бесед позволяет не смешивать canonical model state с текстом сообщений.

\begin{table}[H]
    \centering
    \caption{Ответственность модулей}
    \label{tab:module-responsibility}
    \small
    \begin{tabularx}{\textwidth}{L{3cm}L{6cm}X}
        \toprule
        Модуль & Ответственность & Не входит в его задачи \\
        \midrule
        generation\_core & IR, builders, validation, simulation, formats, benchmark & Учётные записи и разграничение доступа \\
        backend & Аутентификация, беседы, версии, persistence, jobs & Самостоятельное изменение уравнений \\
        frontend & Чат, графы, LaTeX, симуляция, экспорт & Хранение второй версии модели \\
        \bottomrule
    \end{tabularx}
\end{table}

\section{Контекстная работа с моделями}

Вход пользователя обрабатывается как запрос на ответ, уточнение, создание новой модели или изменение существующей. Для новой модели используется обычный конвейер генерации. Для изменения в \codeid{GenerationRequest} передаётся \codeid{base_model}, который служит источником истины, а текстовая история рассматривается только как контекст. Это исключает неявное применение всех предыдущих сообщений как одновременных изменений.

Редактирование требует полного обновлённого документа, а не приблизительного patch. Переменные, уравнения и настройки, не затронутые запросом, должны сохраняться. Каждая новая модель получает собственную версию, поэтому сообщения пользователя и ответы системы отображаются в правильном порядке.

Возможны три типа ответа: текстовый ответ, уточняющий вопрос и модель. Предполагаемое допущение сохраняется с источником и признаком `requires\_confirmation`. Политика `require\_explicit` возвращает `MISSING\_PARAMETER`, если существенные значения требуют подтверждения. Это устраняет прежнюю практику скрытой подстановки ненулевых defaults.

\section{Вычисления и анализ}

REST API поддерживает создание моделей, контекстные беседы, импорт Vensim/XMILE, симуляции, сценарии, анализ чувствительности и калибровку. Операции, требующие длительных вычислений, оформляются как jobs с проверкой состояния и idempotency key. Повторный запуск failed job выполняется через отдельный retry endpoint.

Пользователь может переопределять параметры и начальные значения уровней, выбирать solver и задавать временную сетку. Внешние ряды разрешены только для exogenous variables и должны покрывать весь горизонт. Сценарии могут выполняться с сохранением частичных результатов, чтобы ошибка одного набора параметров не скрывала остальные.

Seeded sensitivity analysis повторяется по конфигурации и seed, а результаты содержат успешные и неуспешные запуски. Калибровка сопоставляет наблюдениям переменные и временные точки, подбирает параметры в заданных границах и создаёт новую версию документа, не изменяя исходную модель. Эти возможности используют один canonical IR, поэтому не создают отдельную копию уравнений для каждого анализа.

\section{Пользовательский интерфейс}

Frontend реализован на Next.js \cite{nextjs2024} и предоставляет:

\begin{itemize}
    \item вход по email/OTP без гостевого доступа;
    \item список бесед и создание новой модели;
    \item контекстные вопросы и последовательные изменения;
    \item CLD/SFD, строящиеся из проверенных связей;
    \item таблицы переменных и уравнения в LaTeX-форме;
    \item выбор переменной и отображение временных рядов;
    \item скачивание JSON, Python, Vensim, XMILE и CSV.
\end{itemize}

Формулы отображаются через KaTeX, а не как необработанный DSL. CLD/SFD не генерируются как независимые картинки: их связи выводятся из документа и сверяются с зависимостями. Это уменьшает вероятность расхождения между визуальной и вычислительной версиями.

\section{Локальные и удалённые провайдеры}

Единый LLM-интерфейс поддерживает fixture, OpenRouter и локальные OpenAI-compatible серверы. Fixture нужен для воспроизводимых unit-тестов, а не для утверждения качества реальной модели. OpenAI-compatible клиент поддерживает ограниченные повторы при сетевых сбоях, strict JSON Schema и диагностику пустого или обрезанного ответа.

Локальный vLLM использует PagedAttention \cite{kwon2023pagedattention}. Для Qwen3-8B-GGUF применяется файл Q4\_K\_M и базовый tokenizer Qwen3-8B \cite{qwen2025technical}. В vLLM 0.30 GGUF требуется отдельный официальный plugin, помеченный как experimental; он не включается в зависимости generation core. Для solver наличие GPU не требуется.

Полная конфигурация локального inference использует контекст 16384 токена, `max\_tokens=8192` и `concurrency=1`. Увеличение `gpu-memory-utilization` приводит к резервированию дополнительной памяти под KV-cache и не должно ошибочно интерпретироваться как размер модели.

Для малых LLM предусмотрены два независимых opt-in режима: полная ComponentPlan JSON Schema и pattern shortcut. Второй режим использует validated archetype без LLM только для качественного описания, все числа которого уже принадлежат эталонному шаблону; пользовательские значения отключают его. По умолчанию оба режима выключены и записываются в метаданные эксперимента. Для явных самолинейных ODE шаг интегрирования ограничивается с учётом скорости и горизонта, а для delay --- отдельно для сохранения третьего порядка. Внутренние шаги ограничены 10000, пользовательский шаг не перезаписывается.

\section{Безопасность и ограничения исполнения}

Платформа не выполняет пользовательский Python, pickle, notebook или сторонний native-код. Все выражения проходят ограниченный AST, а Python создаётся доверенным кодогенератором. Vensim macros разворачиваются bounded parser с ограничением глубины, числа определений и размера результата.

XML DTD и внешние сущности запрещены. Загружаемые модели ограничиваются размером и проверяются по схеме. Work directories изолированы, watchdog обрабатывает зависшие jobs, а конфигурация задаёт пределы памяти, файлов, времени и concurrency. Секреты не включаются в benchmark artifacts.

\section{Проверка реализации}

На момент фиксации текущего черновика полный offline-suite включает 305 тестов generation core, 20 backend-тестов, 12 frontend-тестов и 2 браузерных E2E-сценария. Дополнительно выполняются Ruff, Biome, TypeScript, lockfile checks и production frontend build.

Тесты охватывают нулевые начальные значения, units, topology, repair, explicit ODE, массивы, `:EXCEPT:`, macros, stock-free functions, сценарии, sensitivity, calibration, REST, OTP, отображение уравнений и полный пользовательский сценарий. Они подтверждают инженерную работоспособность покрытых путей, но не заменяют независимую научную оценку генеративного качества.
